\documentclass[11pt]{article}
\newcommand{\LabNumber}{6}
\newcommand{\LabTitle}{Queues and Deques}
\newcommand{\LabTopic}{Queue ADT, circular array, double-ended queue}
\input{common.tex}

\begin{document}
\labheader

\section*{Objectives}
\begin{itemize}[leftmargin=*,nosep]
  \item Define the Queue ADT (FIFO) and the Deque ADT (double-ended).
  \item Implement a queue backed by a circular array using modular arithmetic.
  \item Implement a Deque using a doubly linked list.
  \item Apply queues to a round-robin task scheduler simulation.
\end{itemize}

\begin{summary}[Quick Reference]
\textbf{Queue interface.}
\begin{lstlisting}
void enqueue(E e);  // add to back
E dequeue();        // remove from front
E first();          // peek front, no removal
int size(); boolean isEmpty();
\end{lstlisting}

\textbf{Circular array.}\quad Maintain \texttt{front} and \texttt{size}.
Back index: \texttt{(front + size) \% capacity}.
\begin{lstlisting}
// enqueue:
data[(front + size) % capacity] = e;  size++;
// dequeue:
E val = data[front];
front = (front + 1) % capacity;  size--;
return val;
\end{lstlisting}

\textbf{Deque API.}\quad
\texttt{addFirst}, \texttt{addLast}, \texttt{removeFirst}, \texttt{removeLast},
\texttt{peekFirst}, \texttt{peekLast} --- all $O(1)$.
\end{summary}

\subsection*{Quick Experiments (Try It)}

\begin{tryit}{Wrap-around with a small capacity}
Create an \texttt{ArrayQueue<Integer>} of capacity 4.
Enqueue 1, 2, 3, 4. Dequeue twice. Enqueue 5, 6.
Print \texttt{front} and the logical contents.
Confirm \texttt{front} has wrapped around.
\end{tryit}

\begin{tryit}{Queue vs.\ Stack order}
Push numbers 1--5 onto a stack, then pop them into a queue.
Dequeue and print. What order do you see?
What did the stack-to-queue transfer accomplish?
\end{tryit}

\section*{In-Lab Exercises (target: 2 hours)}

\begin{labexercise}{Array-Based Queue \hfill {\small \itshape (40 min)}}
Implement \texttt{ArrayQueue<E>} backed by a fixed-capacity \texttt{Object[]}.
\begin{itemize}[leftmargin=*,nosep]
  \item Fields: \texttt{data[]}, \texttt{front}, \texttt{size}, \texttt{capacity}.
  \item \texttt{enqueue} throws \texttt{IllegalStateException} when full.
  \item \texttt{dequeue} and \texttt{first} throw \texttt{NoSuchElementException} when empty.
  \item \texttt{toString()} shows front-to-back, e.g.\ \texttt{front->[1, 2, 3]<-back}.
\end{itemize}
Test with wrap-around: capacity 5, enqueue 5, dequeue 2, enqueue 2 more.
\end{labexercise}

\begin{labexercise}{Linked-List Deque \hfill {\small \itshape (30 min)}}
Using \texttt{DoublyLinkedList<E>} from Lab~4, wrap it in \texttt{LinkedDeque<E>}
with the full Deque API. All operations must be $O(1)$.
Write \texttt{boolean isPalindrome(String s)} using the deque.
Test on \texttt{"level"}, \texttt{"world"}, \texttt{"noon"}.
\end{labexercise}

\begin{labexercise}{Round-Robin Scheduler \hfill {\small \itshape (30 min)}}
Simulate a round-robin CPU scheduler:
\begin{enumerate}[leftmargin=*,nosep]
  \item \texttt{Task} has fields \texttt{name} (String) and \texttt{burstTime} (int ms).
  \item Read at least 4 tasks from the user.
  \item Time quantum = 3~ms. Each round: dequeue, subtract quantum (or finish if $\leq$ quantum),
        print a log line, re-enqueue if remaining time $> 0$.
  \item Continue until the queue is empty.
\end{enumerate}
Example log: \texttt{[t=6] Running TaskA, remaining=4ms}.
\end{labexercise}

\begin{aicollab}
\textbf{Suggested prompts:}
\begin{itemize}[leftmargin=*,nosep]
  \item ``Why does a circular array avoid wasting space that a naive queue wastes?''
  \item ``What is the difference between \texttt{Deque} and \texttt{Queue} in Java's standard library?''
\end{itemize}
\medskip\textbf{Verify:} Ask the AI to draw the circular array state after the wrap-around sequence
in Try It 1. Compare with your own diagram.
\end{aicollab}

\takehomebanner

\begin{trace}[Circular array trace]
Trace on an \texttt{ArrayQueue} of capacity 4.
After each operation record \texttt{front}, \texttt{size}, and the array (\texttt{-} for empty slots).
\begin{lstlisting}
enqueue(A); enqueue(B); enqueue(C); dequeue();
enqueue(D); enqueue(E); dequeue(); dequeue();
\end{lstlisting}
\end{trace}

\begin{writecode}[Reverse a queue using a stack]
Write \texttt{void reverseQueue(Queue<Integer> q)} that reverses all elements in \texttt{q}
using only one auxiliary \texttt{Stack<Integer>}. No other data structures allowed.
\end{writecode}

\begin{drawex}{Round-robin Gantt chart}
Draw a Gantt chart for these tasks with quantum = 2~ms:
\medskip
\begin{tabular}{@{}lll@{}}
\toprule
Task & Burst & Arrival \\ \midrule
P1 & 6 ms & 0 \\ P2 & 4 ms & 0 \\ P3 & 2 ms & 0 \\ P4 & 5 ms & 0 \\
\bottomrule
\end{tabular}
\end{drawex}

\vspace{4pt}
\noindent\textbf{Submission.} Save files as \texttt{lab6\_ex*.java} and submit on the course site.

\end{document}
